\chapter{知识演进与迁移的评价}
\label{chap:ket-evaluation}

花卉识别系统进入新地区后，整体准确率上升，稀有病害的漏检却增多了。经过第三个季节的更新，旧地区的成绩高于最初上线时，又低于上一版本。物种模型与病害模型合并后，各自的问答仍能完成，但系统找不出“指定物种中患病的花朵”。我们应当怎样判断这些改变？整体正确数、相对历史版本的变化和跨任务组合能力分别回答不同问题，其中一项改善不能替另一项作结论。

第~\ref{chap:ket-reuse-selection}~章已经区分固定预算下的效果与达到同一效果的代价；第~\ref{chap:ket-update-maintenance}~至第~\ref{chap:ket-sharing}~章给出了改变模型或系统的机制。这里接着问：需要留下什么记录，才能检验这些机制是否实现了预期目标？只有先明确比较对象和指标，才能判断数据是否覆盖了所声称的变化；只有把信息权限和对照条件固定下来，结果差异才有可追查的解释。最后还要考察收益落在哪些任务上、付出了什么资源，以及证据尚未涉及哪些条件。本轮失败可以成为下一轮研究的起点，却不能用来事后重定本轮的成功标准。

收益方向仍采用式~\eqref{eq:ket-reuse-risk-gain}的约定：$R_T$是目标分布上的期望损失，$f_0$不使用待检验来源，$f_K$使用该来源，$\Delta_T=R_T(f_0)-R_T(f_K)$。损失下降给出正收益；对准确率这样的增大型指标，则用方法分数减去对照分数。下文的花卉照片、三任务成绩、机构事实和费用都是人工教学构造，用于核对定义与计算，并非实际测量。初始学习曲线及来源准备费取自第~\ref{chap:ket-reuse-selection}~章；引用的基准研究则用于说明真实评价协议与其边界。

\section{评价目标与指标的确定}
\label{sec:ket-evaluation-metrics}

假如上线方案只允许使用80个适应标签，评价要比较这个预算下的新地区错误率；假如要求达到90\%准确率，评价就要追踪各方案为此使用的全部开发标签。这是两项可分别被数据否定的主张。同样，要求旧地区仍能正确识别花卉，不等于要求模型保留所有过时答案。因而每项主张都应写明被评价的对象、对照、预算和达标条件。多个目标可以按事先约定的权重汇总，但原始结果必须保留，否则无法判断关键任务是否被其他任务的改善掩盖。

\subsection{目标表现与适应效率的度量}
\label{sec:ket-evaluation-target}

要回答固定预算下哪个方案更合适，就让两个方案在相同的预定预算点接受测量。以$s_r(b)$记录第$r$次运行在预算$b$处的目标分数；横轴可以采用标签数、更新步数或设备时间，但必须保持同一口径，不能把各方案最有利的检查点拿来相减。每次运行保存数据划分、种子、来源版本、可见目标信息、分数的计数依据及累计资源。

在同一目标样本和可匹配随机条件下运行两个方案，可以先求每对运行的收益$d_r(b)$：增大型指标取方法减对照，损失取对照减方法。这保留了两个结果来自同一条件的对应关系。要同时描述平均收益和运行间波动，对$R$个配对差求均值及样本标准差：
\begin{equation}
  \bar d(b)=\frac1R\sum_{r=1}^{R}d_r(b),\qquad
  s_d(b)=\sqrt{\frac{1}{R-1}
    \sum_{r=1}^{R}\bigl(d_r(b)-\bar d(b)\bigr)^2}.
  \label{eq:ket-evaluation-run-summary}
\end{equation}
这里要求运行数$R>1$；$s_d$描述配对收益在运行间的离散程度，并不直接给出均值的置信区间。同一批困难样本可能同时降低两个方案的成绩，配对差仍保留这种联系。若分别挑出两个方案的最优运行，差值就不再来自同一条件。用于调参的验证记录与锁定方案后的独立测试记录也应分开。

一次配对测量的教学数据见图~\ref{fig:ket-reuse-learning-curves}。四个适应标签预算点组成集合
\[
  N=\{20,40,80,160\}.
\]
在这些预算上，两种方法的准确率为
\[
  \begin{aligned}
    \text{迁移} &: (82,88,91,92)\%,\\
    \text{从头训练} &: (65,76,85,91)\%.
  \end{aligned}
\]
两列均按同一组100张验证照片计算。逐点相减得到$(17,12,6,1)$个百分点：在这些已测点上，迁移的优势随预算增大而缩小。横轴没有计入来源筛选与验证标签；若要比较全部开发标签，应按各方案实际用途重新核算横坐标，而不是仅更换轴名。

再构造两次运行，仍使用相同预算点和同一组100张验证照片。第二次迁移正确数为$(80,87,90,91)$，从头训练为$(64,75,84,90)$；第三次对应为$(81,86,89,89)$与$(66,77,85,89)$。在80标签处，配对收益为$(6,6,4)$个百分点，式~\eqref{eq:ket-evaluation-run-summary}~给出均值$16/3\approx5.33$和样本标准差$\sqrt{4/3}\approx1.15$个百分点。第一条曲线的91\%不能代替这三次记录；300次预测也不是300张独立照片，因为每张照片被重复评价了三次。

若上线要求是达到指定效果，就需要从曲线中找出首次满足要求的已测预算。对增大型分数，先固定阈值$s_\star$和观测预算集合$B$，再定义第$r$次运行的经验到达量：
\begin{equation}
  b_{\star,r}=\min\{b\in B:s_r(b)\geq s_\star\}.
  \label{eq:ket-evaluation-threshold}
\end{equation}
最小值只在集合非空时存在。若所有已测点都未过线，记“截至最大预算$b_{\max}$未达到”，不把终点填成成功成本。把继续按既定测量规则观察的首次过线视为待观察事件时，这属于\term{右删失}（right censoring）：观察已经结束，事件尚未被记录。它不说明未测中间预算的成绩，也不保证以后会过线。

式~\eqref{eq:ket-evaluation-threshold}~采用单点过线规则；若要求连续多个点达标或置信下界过线，应在测量前相应改定判据。依据验证成绩停止训练后，还需在独立测试上评价选定模型。经验阈值只决定本次运行在哪个已测点停止，不能保证总体准确率。

以90\%为阈值，上述三次迁移记录为$(80,80,>160)$，从头训练为$(160,160,>160)$，其中$>160$表示截至观测终点仍未记录到过线。两方案截至160的到达比例都是$2/3$，但截至80分别为$2/3$和0。这些比例保留了未成功的第三次运行；只对成功运行求条件均值，则得到80与160，不能称为全部运行的平均成本。

也可以事先规定把未到达者按观察上限截断，计算迁移的$(80+80+160)/3\approx106.67$和从头训练的160。此时应称为截断预算统计量，第三个160只是停止观察时已经允许的预算。它既不是未删失的平均到达量，也不证明第三次运行最终能够成功。

观测间隔限制了到达量的精度。第一条曲线只能确认迁移在已测点80过线、从头训练在160过线；40与80之间没有测量，插值不能产生精确标签成本。计算曲线面积也需指定横轴是标签数$n$还是$\log n$、积分区间及插值方法；等比预算点上的分数均值并不等于按标签投入积分。跨任务汇总还会遇到单位差异：准确率、像素误差和推理成功率应先各自确定收益方向，再依据预定尺度与权重比较。

零样本成绩同样依赖允许的信息。类别名称、属性描述和预训练知识可以参与预测，但若已经用目标标签挑选提示，就应报告这份监督，不能只因预测时没有更新参数而称其完全无目标监督。少样本评价需列明类别数、每类支持样本数、查询集，以及额外可见的无标签输入。

当新物种还不在原输出空间中，评价要区分发现未知、学习新类与校准旧新类别三项要求。未知检测同时报告已知类的误拒率和未知类的误接纳率；新类学习在新类留出照片上计分；校准检查旧新类别共同竞争时的结果。全部拒绝陌生照片只能提供某种检测行为，不能证明系统已经会识别新增物种。

\subsection{持续能力变化的度量}
\label{sec:ket-evaluation-continuous}

要判断第三次更新是否损害旧地区，需要知道同一任务在不同时刻的成绩，而不只是最终总分。设系统依次学习$K$个任务，阶段$t$结束得到$f^{(t)}$，在任务$k$固定评测集$Q_k$上的准确率为$a_k^{(t)}$。矩阵$\bm A$以阶段为行、任务为列，元素$\bm A_{t,k}=a_k^{(t)}$；第0行记录尚未经历该序列的共同起点。沿同一列比较才能追踪同一能力。

对角线$a_k^{(k)}$记录任务刚学完时的成绩；$k<t$与$k>t$分别对应历史任务和尚未学习的任务。后者仍可能借助已有输出接口做零样本预测。没有测量或缺少可用接口时，应保留缺测，而不是填入零分。

GEM采用这样的阶段与任务记录，分别汇总最终准确率、相对对角线的后向变化和相对随机初始化基线的学前变化。本章用$a$表示成绩，以免与风险$R$混淆。若两个方案共享预训练起点，学前对照应采用尚未经历历史序列的该起点，并说明这已不同于GEM的随机初始化基线；比较的仍是额外历史经验的作用\scite{ket-lopezpaz2017-gem}

\subsubsection{历史能力保持与后向迁移的度量}
\label{sec:ket-evaluation-backward}

假设系统在阶段$t$已经承担前$t$项任务，先问当前模型在这些任务上的整体成绩怎样。任务等权的\term{阶段平均表现}（stage-wise average performance）把当前行中已经学习的任务准确率求平均：
\begin{equation}
  \bar a^{(t)}=\frac1t\sum_{k=1}^{t}a_k^{(t)}.
  \label{eq:ket-evaluation-stage-average}
\end{equation}
分母$t$计的是已学习任务，不是计划中的全部$K$项，也不是照片数。若任务重要性不同，可以使用预定权重，但要说明随着任务到达是否重新归一化；包含未学任务的平均应另报。由于平均范围随$t$扩大，新加入的困难任务也会拉低$\bar a^{(t)}$。因此，这条曲线不能单独回答旧能力是否下降。

旧能力的变化要固定任务列，再选参照时刻。对历史任务$k<t$，\term{遗忘量}（forgetting）$F_k^{(t)}$问当前成绩比过去已测峰值低多少，所以先在任务学完至当前阶段之前的检查点中取最大值。\term{后向迁移}（backward transfer，BWT）$B_k^{(t)}$则问后续学习使该任务比刚学完时改变多少，参照是对角线。两种逐任务差值及历史任务等权平均定义为：
\begin{equation}
  \begin{aligned}
    F_k^{(t)}&=\max_{k\leq u<t}a_k^{(u)}-a_k^{(t)},\\
    \bar F^{(t)}&=\frac1{t-1}\sum_{k=1}^{t-1}F_k^{(t)},\\
    B_k^{(t)}&=a_k^{(t)}-a_k^{(k)},\\
    \operatorname{BWT}^{(t)}&=\frac1{t-1}\sum_{k=1}^{t-1}B_k^{(t)}.
  \end{aligned}
  \label{eq:ket-evaluation-forgetting-bwt}
\end{equation}
索引$u$只遍历已保存的历史检查点，不含当前阶段。$F>0$表示低于历史峰值，$B>0$表示超过刚学完的成绩；当前达到新高时，$F$可以为负。若只关心损失，可以另报截断量$\max(0,F)$，但不能把它默认为原定义。阶段1尚无历史任务，两个平均的分母为零，应记为未定义而非零。历史检查点越密，越可能观测到较高峰值，故跨方案比较遗忘量还需匹配测量频率。

下面把同一个两物种识别要求放入三个依次到达的环境：$\mathcal T_1$为地区甲春季，$\mathcal T_2$为地区甲夏季，$\mathcal T_3$为地区乙秋季。此处编号表示学习序列，不沿用第~\ref{chap:ket-sharing}~章的物种、病害和定位三项任务。标签含义与输出空间不变，模型不接收任务身份；每个任务另留100张固定照片，不进入训练。

图~\ref{fig:ket-evaluation-matrix}~给出人工正确数，每列分母均为100，因此格内数值也等于准确率百分数。共同起点在三项任务上都是50\%。阶段1没有评价任务3，该格保留缺测；右侧同列比较指出两个历史参照怎样产生不同答案。

\begin{figure*}[htbp]
  \centering
  % Teaching data, not measurements. Every task has 100 fixed evaluation images.
% Rows: (50,50,50), (80,55,missing), (86,82,60), (83,78,88).
\begingroup
\setlength{\tabcolsep}{2.5pt}
\setlength{\fboxrule}{0.6pt}
\renewcommand{\arraystretch}{1.65}
\begin{tikzpicture}[
  x=1mm,y=1mm,
  font=\figurefont\small,
  text=BookInk,
  cell/.style={minimum width=12mm,minimum height=7mm,inner sep=1mm,
    draw=BookMuted,line width=0.6pt},
  flow/.style={-{Latex[length=2mm]},draw=BookMuted,line width=0.8pt}
]
  \node[anchor=north west,inner sep=0pt,font=\figurefont\small\bfseries]
    at (0,4) {(a) 阶段 $\times$ 任务};
  \node[anchor=north west,inner sep=0pt] at (0,-3) {
    \begin{tabular}{@{}p{26mm}*{3}{>{\centering\arraybackslash}p{11mm}}@{}}
      \toprule
      阶段 & $\mathcal T_1$ & $\mathcal T_2$ & $\mathcal T_3$ \\
      \midrule
      0：共同起点 & 50 & 50 & 50 \\
      1：学完$\mathcal T_1$ &
        \fcolorbox{BookTeal}{BookTeal!10}{80} &
        \colorbox{BookAmber!12}{55} &
        \colorbox{BookAmber!12}{\textemdash} \\
      2：学完$\mathcal T_2$ &
        \colorbox{BookBlue!12}{86} &
        \fcolorbox{BookTeal}{BookTeal!10}{82} &
        \colorbox{BookAmber!12}{60} \\
      3：学完$\mathcal T_3$ &
        \colorbox{BookBlue!12}{83} &
        \colorbox{BookBlue!12}{78} &
        \fcolorbox{BookTeal}{BookTeal!10}{88} \\
      \bottomrule
    \end{tabular}
  };
  \node[anchor=west,inner sep=0pt,font=\figurefont\footnotesize]
    at (0,-49) {$k<t$：历史\quad $k=t$：当前\quad $k>t$：未学};
  \node[anchor=west,inner sep=0pt,font=\figurefont\footnotesize]
    at (0,-56) {单位：准确率（\%）；\textemdash：未测};

  \node[anchor=north west,inner sep=0pt,font=\figurefont\small\bfseries]
    at (81,4) {(b) 同列、不同时刻的比较};
  \node[anchor=west,inner sep=0pt] at (81,-6)
    {遗忘：任务1，历史最佳减当前};
  \node[cell,fill=BookBlue!12] (peak) at (88,-16) {86};
  \node[cell,fill=BookBlue!12] (nowf) at (120,-16) {83};
  \draw[flow] (peak.east) -- (nowf.west);
  \node[anchor=west,inner sep=0pt] at (131,-16) {$F_1^{(3)}=3$};

  \node[anchor=west,inner sep=0pt] at (81,-28)
    {BWT：任务1，当前减刚学完};
  \node[cell,draw=BookTeal,fill=BookTeal!10] (learned) at (88,-38) {80};
  \node[cell,fill=BookBlue!12] (nowb) at (120,-38) {83};
  \draw[flow] (learned.east) -- (nowb.west);
  \node[anchor=west,inner sep=0pt] at (131,-38) {$B_1^{(3)}=3$};

  \node[anchor=west,inner sep=0pt] at (81,-50)
    {零样本FWT：任务2，学前减起点};
  \node[cell,fill=BookPaper] (base) at (88,-60) {50};
  \node[cell,fill=BookAmber!12] (before) at (120,-60) {55};
  \draw[flow] (base.east) -- (before.west);
  \node[anchor=west,inner sep=0pt] at (131,-60) {$Z_2=5$};
\end{tikzpicture}
\endgroup

  \caption{三任务持续学习的教学矩阵。行是完成的阶段，列是评测任务；每格按100张固定照片计算准确率（\%）。实线框标出刚学完的对角线，浅底区分历史与未学区域，不编码成绩高低；\textemdash 表示未测。右侧保持任务列不变，分别以历史峰值、刚学完成绩和阶段0起点为参照计算遗忘、BWT与零样本FWT。}
  \label{fig:ket-evaluation-matrix}
\end{figure*}

按已学任务平均，三个阶段得到$\bar a^{(1)}=80\%$、$\bar a^{(2)}=(86+82)/2=84\%$与$\bar a^{(3)}=(83+78+88)/3=83\%$。阶段2只有任务1属于历史任务：其成绩由80\%升至86\%，故$F_1^{(2)}=80-86=-6$个百分点，$B_1^{(2)}=86-80=6$个百分点。这里负遗忘与正后向变化都记录了成绩创新高；它们本身没有解释改善机制。

\begin{samepage}
阶段3要比较两个历史任务。任务1的历史峰值86\%降至83\%，因此$F_1^{(3)}=\max(80,86)-83=3$个百分点；任务2从刚学完的82\%降至78\%，因此$F_2^{(3)}=82-78=4$个百分点。平均遗忘为$(3+4)/2=3.5$个百分点，分母只包括这两个任务。若把参照改为各自刚学完的成绩，得到：
\[
  \operatorname{BWT}^{(3)}
  =\frac{(83-80)+(78-82)}2
  =-0.5\text{个百分点}.
\]

\end{samepage}
任务1比峰值低3个百分点，又比刚学完高3个百分点，因而“有所遗忘”和“后向获益”可以同时成立。把两种定义相加，当前成绩恰好抵消，得到：
\[
  F_k^{(t)}+B_k^{(t)}
  =\max_{k\leq u<t}a_k^{(u)}-a_k^{(k)}\geq0.
\]
最大值的候选中包含刚学完时刻$u=k$，所以右侧非负。只有对角线成绩等于此前峰值时，两种差值才满足$F_k^{(t)}=-B_k^{(t)}$；否则它们不能互相替代。

逐列求差还要求比较的是同一有效目标。例~\ref{ex:ket-overview-dated-fact}与第~\ref{sec:ket-update-edit-validation}~节中的甲所于2026年7月1日由甲交接给乙。评价历史任职时仍应答甲，评价交接后的现任时应答乙；纠正后一类问题的过时答案不是有害遗忘。问题及标准答案都需带时间条件，历史、当前和未变任务分别计分。物种分类标准发生变化时，也应保留标准版本或建立明确映射，不能直接相减语义不同的成绩。

同一批旧照片还可能通过不同接口接受测试。提供任务身份以选择旧任务头，评价的是带路由的多头系统；不给身份而让全部类别竞争，增加了统一输出空间的判别要求。若允许调用历史检查点，保持成绩则属于模型与版本库组成的系统，不再只属于当前模型。报告应据此列出任务身份权限、实际路由正确率、检查点访问范围和存储费用。

\subsubsection{前向迁移与持续可学性的度量}
\label{sec:ket-evaluation-forward}

同一矩阵还能检验历史经验是否帮助后来的任务，但要把学习前的预测与学习后的效率分开。旧任务保持良好并没有回答其中任何一项。\term{零样本前向迁移}（zero-shot forward transfer）比较新任务尚未接受训练时的成绩，问先前经验能否直接改善预测。以共同起点成绩$b_k=a_k^{(0)}$为参照，将任务$k$紧邻学习前的成绩减去它，再平均任务2至$K$的差值：
\begin{equation}
  Z_k=a_k^{(k-1)}-b_k,\qquad
  \operatorname{FWT}_0
  =\frac1{K-1}\sum_{k=2}^{K}Z_k.
  \label{eq:ket-evaluation-zero-forward}
\end{equation}
任务1之前没有本序列的经验，故分母为$K-1$，并要求$K>1$。$Z_k>0$表示尚未学习任务$k$时已优于共同起点。这个比较需要相同的标签含义、类别描述和预测接口；若必须用新任务标签拟合输出头才能预测，就没有该接口下的零样本成绩，应记为不适用，不能用拟合后的分数补入。

图~\ref{fig:ket-evaluation-matrix}~中，任务2与任务3紧邻学习前分别为55\%和60\%，因此$Z_2=55-50=5$、$Z_3=60-50=10$个百分点，平均为$\operatorname{FWT}_0=(5+10)/2=7.5$个百分点。公式不需要阶段1的任务3成绩，故该处缺测不妨碍计算。若改问任务1的学习单独怎样影响任务3，就需要补测该格，不能由阶段2的60\%还原它。

CLEAR补充材料中的同名指标使用了不同参照：Backward Transfer平均严格下三角的绝对准确率，Forward Transfer平均严格上三角的绝对准确率，均未求相对基线的差。本例下三角为$(86+83+78)/3\approx82.33\%$，与本章$-0.5$个百分点的BWT不是同一个量；上三角缺少$a_3^{(1)}$，不能完整求平均。其Next-Domain Accuracy只取紧邻上对角线，本例为$(55+60)/2=57.5\%$，而非本章减去50\%起点后的7.5个百分点。使用文献名称时必须同时给出所取矩阵区域、运算和单位\scite{ket-lin2021-clear}

学前成绩只回答能否直接预测，学习效率还需比较新任务的适应曲线。设$c_{\star,k}^{\mathrm{hist}}$和$c_{\star,k}^{0}$分别为带历史经验的模型与匹配无历史对照在任务$k$达到同一预定阈值的适应成本。若两者实际过线且对照成本为正，用节省量除以对照成本，得到相对改善：
\begin{equation}
  E_k=\frac{c_{\star,k}^{0}-c_{\star,k}^{\mathrm{hist}}}
             {c_{\star,k}^{0}}.
  \label{eq:ket-evaluation-learning-cost-gain}
\end{equation}
$E_k>0$表示省成本，$E_k<0$表示额外花费；归一化参照是$c_{\star,k}^{0}$，不是两方案的成本之和。为与矩阵对应，构造预算点$n=(20,40,80,160)$：任务2的历史模型与无历史对照成绩分别为$(82,90,92,93)\%$和$(70,82,90,92)\%$，任务3分别为$(76,88,90,92)\%$和$(65,76,85,90)\%$。各预算点采用嵌套支持集、同一每类抽样规则及100张固定评测照片，并匹配容量与优化预算。90\%阈值下，任务2的到达量为40与80，任务3为80与160，所以$E_2=(80-40)/80=50\%$、$E_3=(160-80)/160=50\%$。

原序列学习任务2时只用了20个标签，学习任务3时用了40个，因此阶段终点仍为图中的82\%与88\%。扩展曲线由阶段1、2模型的独立副本产生，只在隔离探测中追加预算，探测模型及反馈不返回原序列。这样测到的是同一起点的后续学习效率，而不是借额外评价训练改变原系统的历史。探测花费仍须登记为评价成本。

成本节省50\%与学前准确率增益7.5个百分点采用不同单位，不能直接合成“前向迁移分数”。若任务3截至160标签仍未过线，应按前述规则报告$c_{\star,3}^{\mathrm{hist}}>160$及到达比例，不代入一个虚构的有限成本。模型在零成本时已达标，可以直接报告零样本成功；若对照成本也为零，式~\eqref{eq:ket-evaluation-learning-cost-gain}~不能作零除。

一次新任务学习较快，还不足以判断多次更新后是否仍能吸收新知识。\term{持续可学性}（continual learnability）考察这种后续学习能力，应从多个阶段的副本出发，学习同一个隔离探测任务。固定新样本、输出头、容量、调参预算和停止规则，再比较曲线，仍不将探测反馈交回原序列。

设三个历史阶段的旧任务成绩都为90\%，同一探测任务的过线步数却依次为40、80、160，匹配的无历史对照始终需要80步。相对成本改善便为$50\%,0,-100\%$。这个教学反例说明保持旧分数并不保证以后容易学习；它支持后续学习能力下降的行为描述，下降原因还需用新的对照检验。

若后期模型扩容，探测应增加同容量的无历史对照，并列出新增存储和计算，避免把容量增长全归给历史经验。任务排列变化要通过顺序重复检验，只换初始化无法覆盖它。长期主张还需要包含旧环境再次到达和不同难度任务，并将历史训练与当前探测分别记账。单次适应成功的证据范围仍限于那一次学习。

\subsection{知识修正与关联一致性的度量}
\label{sec:ket-evaluation-editing}

假设甲所迁至另一座城市，直接问所在地与问出差应去哪一省，就需要检验事实及其不同使用后果。知识编辑评价因此要分开四组问题：原请求检验修改是否完成，未参与编辑的等义改写或跨语言表达检验表达泛化，依赖新事实的推理题检验后果传播，语义上无关的问题检验保持。各组分别报告成功率，并注明按请求还是按问题计分。若一条请求配十种改写、另一条只有一种，先在请求内求正确率再等权汇总，才不会让前者获得十倍权重。

这些成功率还依赖单题如何判分，应在运行前冻结规则。封闭答案可按约定完成Unicode规范化、大小写及空白处理，再匹配人工确认的等价答案；日期、数值与单位使用指定解析器和容差。开放生成若由模型评分，须固定评分器版本、提示、证据权限和输出标签，并在不参与方法选择的人工标注子集上报告评分一致性。无法解析、拒答及多答案冲突也需预定处理。单独命中关键词不能证明语义正确，事后见到有利答案也不能回写判分标准。

为具体区分四组问题，给甲所增加一项独立于负责人交接的迁址请求：自2026年9月1日起从甲城迁至乙城，两城始终分别属于甲省与乙省。“9月以后所在地是哪座城市”及“9月应到哪座城市拜访”都应答乙城；“9月出差应前往哪个省”则需将新所在地接入“城市所属省”的关系，得到乙省。问2026年8月应去哪个省仍应答甲省；另一家未迁址机构的地址不应改变。时间与依赖关系由构造明确给出，不靠模型猜测。

构造10条编辑请求，每条配1道直接题、2道未用于编辑的等义改写和1道关联推理题。直接题答对9道、改写答对16道、推理答对6道时，请求内平均后再等权汇总，分别得到90\%、80\%与60\%；由于每条请求的题数相同，这里恰好也等于各组按题统计的比例。直接成功并未覆盖关联推理失败。

再设100道无关题，98道输出未变，输出保持率为98\%。它不等于正确率：原模型的错误也可能被保留。无关集应同时报告编辑前后准确率及预测变化率；若还检验概率分布变化，要预先定义距离和判断阈值。这样才能区分保护有效内容与仅维持原输出。

迁址导致省份答案随之改变，是修改目标的合理后果；\term{局部性}（locality）要求保护的是不依赖该修改的内容，而非让所有邻近答案不变。MQuAKE用多跳问题检验编辑后的关联后果，并分别设置反事实修改与真实时间更新。其证据限于具有正确前提和已知关系的有限事实链，不能推广为全部推理的一致性保证\scite{ket-zhong2023-mquake}

保持对象还包括不同层次的要求。物种分类成绩未降，不能说明历史负责人问答正确；事实问答正确，也不能说明未授权记录不会被输出。保护集应从修改目标与使用要求确定，不能只选容易保持的通用题。连续编辑后逐次复查仍有效请求、带时间条件的历史请求和相互依赖请求，才能追踪错误是否累积。

如果请求不是修改所在地，而是撤回曾用于训练的指定记录，评价对象就从答案内容转为训练影响。这是机器遗忘要求的比较：遗忘结果应与未使用撤回记录的训练结果对照。以$S$为原训练集、$D\subseteq S$为按记录身份指定的撤回集合，$\mathcal A$为原训练程序、$\mathcal U$为遗忘程序，两个模型写为：
\begin{equation}
  f_{\mathrm{unlearn}}=\mathcal U(\mathcal A(S),D),
  \qquad f_{\mathrm{retrain}}=\mathcal A(S\setminus D).
  \label{eq:ket-evaluation-unlearning-reference}
\end{equation}
随机训练下，严格参照比较训练与遗忘程序产生的发布模型分布，不要求两次参数逐位相等。有限输入上的输出分布可以另作行为检查，却不等价于模型分布保证。参照还要写明初始化、预处理、超参数选择及受撤回记录影响的上游产物；从见过这些记录的检查点继续训练，不是当然合格的重训练。SISA通过限制记录的影响范围、重训受影响部分降低撤回开销，其实验同时报告准确率、重训时间与额外存储\scite{bourtoule2021sisa}

经验检查可在遗忘模型与重训练模型上比较保留任务效用、撤回记录相关行为和预定攻击结果。撤回样本准确率并不必然降为零，因为保留数据仍可能支持相同答案；有限提示拒答或一次攻击失败也不证明训练影响已消失。严格或近似删除还须在规定观察权限下比较相应发布分布，具体标准及验证要求可参考\BookChapterRef{chap:trust-privacy}{隐私保护}。这里记录相应效用、时间和资源，不用普通准确率代替删除要求。

\subsection{多任务共享收益的度量}
\label{sec:ket-evaluation-sharing}

回到第~\ref{chap:ket-sharing}~章共同交付的物种、病害与定位任务，评价要看每项任务相对独立学习改变了多少。以$s_k^{\mathrm{joint}}$和$s_k^{\mathrm{ind}}$分别表示共同与独立学习的同一指标；本节用$d_k=1$表示增大型指标、$d_k=-1$表示减小型指标，不再指配对运行差。要汇总不同单位的收益，还需先约定正尺度$c_k$及任务权重$w_k$，将各任务差值转为：
\begin{equation}
  g_k=d_k\frac{s_k^{\mathrm{joint}}-s_k^{\mathrm{ind}}}{c_k},
  \qquad \bar g=\sum_{k=1}^{K}w_kg_k,\quad
  w_k\geq0,\quad\sum_kw_k=1.
  \label{eq:ket-evaluation-multitask-gain}
\end{equation}
方向系数使$g_k>0$始终表示获益，但除以$c_k$并不会自动产生共同效用。只有事先把$c_k$规定为政策或效用尺度，例如每项任务距最低服务要求允许变化的幅度，才能解释跨任务的一单位收益。用非零独立基线的绝对值作尺度，则只是相对变化指数；基线接近零还会使它不稳定。权重与尺度必须在看结果前固定，原单位结果仍需逐项保留。

设独立方案的物种准确率、病害召回率与定位平均误差为80\%、70\%、10像素，共享方案为84\%、68\%、9像素。教学数据逐项给出$+4$个百分点、$-2$个百分点和误差减少1像素。取独立基线绝对值作$c_k$，相对变化指数为$5\%,-2.86\%,10\%$，等权平均约4.05\%。换一种指标表示或尺度，平均指数也会改变，所以它不是跨任务共同效用。

这组指数的最差值为$-2.86\%$，按点估计受损的任务占$1/3$。它们指出正平均遗漏了病害任务的损失。若病害是主任务，或者共同交付要求不容许它退步，其他任务获益不能抵消违反要求的事实。

最差值要在可比较的尺度上取：80\%与10像素不能直接求最小，可以比较归一化收益，或各自报告距最低服务要求的余量。受损比例也依赖判据。按$g_k<0$的点估计计数，与只将区间完全小于零的任务记为有证据受损，描述不同对象；后者还需单列方向不确定的任务，不能把它们全部记为获益。

若不愿用某项收益补偿另一项损失，可以逐目标比较。\term{帕累托支配}（Pareto dominance）要求A在统一方向的全部目标上不差于B，且至少一项严格更好。上述共享方案因病害召回率下降而不支配独立方案。另一个方案$(84\%,71\%,9\text{像素})$按点估计支配$(84\%,68\%,9\text{像素})$，但此时只比较了三项任务指标；加入计算与存储目标后，还需重新逐项判断。

实际验收也可允许预定幅度的退步，但这种取舍规则不同于不容许退步的帕累托支配。有限样本下应报告各目标配对差及同时覆盖这些目标的区间，再与零或预定容忍边界比较。点估计在边界内，并未排除更大退步；区间跨界时结论仍不确定。训练中的共同下降方向只涉及当时的损失，不能替代最终任务指标的上述检验。

预算参照也需逐项固定。假设三个独立模型各用10 GPU小时，总计30；共享方案让各任务获得规定次数的更新，实际共用24 GPU小时。它可以比较相同任务暴露下的效果，却不能称共享只用了10 GPU小时。若整组任务总预算为10，独立方案必须按预定规则在三模型间分配，共享方案也只能用10，才是同总资源比较。任务暴露次数、标签齐全程度、缺失标签处理、共享骨干费用和调参投入都应列出，更新机会相同不等于计算相同。

\subsection{来源能力保持与组合泛化的度量}
\label{sec:ket-evaluation-composition}

合并物种与病害模型后，能分别回答原问题并不保证能找到指定物种中的患病对象。整合评价因此分开原能力保持、联合请求完成率和未见组合上的正确行为。固定来源$k$的原评测集，以$s_k^{\mathrm{source}}$记录来源模型成绩、$s_k^{\mathrm{merge}}$记录整合系统成绩；沿用上一节方向系数$d_k$，逐来源变化为$d_k(s_k^{\mathrm{merge}}-s_k^{\mathrm{source}})$。参照必须是该来源本身，并在来源等权平均之外列出最受损来源。

例如来源在物种与病害任务上分别为90\%、86\%，融合后为88\%、87\%，逐项变化为$-2,+1$个百分点，等权平均由88\%降至87.5\%。若单任务服务请求中10\%只问物种、90\%只问病害，流量加权成绩则由$0.1\times90+0.9\times86=86.4\%$升至$0.1\times88+0.9\times87=87.1\%$。这个流量平均只汇总两类单任务请求，并不测量同时要求两种能力的联合请求完成率；后者须在同一请求上按完整要求另行计分。等权平均与流量平均各有权重含义，都不能取代物种退化的逐项记录。是否提供任务身份以选择头、是否要求统一标签空间预测，也应分列，避免混合不同难度。

要确认改善超出了来源选择，需要为不同操作设置直接对照。最佳单来源回答全部请求，提供单来源参照；路由器每次选一个来源，检验选择能力；预测集成让多个来源参与，但未必计算新关系；串接模块还要求下游正确使用中间结果。来源调用记录与路由命中率帮助追踪过程，却不能把调用两个模型本身算作组合成功。

花卉例子可通过属性组合留出建立这项检验。训练中让红、黄两色和斑点、条纹两种纹理分别在多个对象与背景上出现，但从来源训练、路由训练、提示示例及组合选择集排除完整的“红色且有斑点”组合。测试要求找出同时满足两属性的对象，分别检查单属性判断与完整合取；同场景加入红色条纹和黄色斑点对象，使只用一个属性的预测无法完成要求。若基础模型预训练已见过相同组合，就只能声称相对本次任务训练留出，不能声称相对全部既有经验未见。

这一划分还需确保合取操作有依据。例如训练中提供其他属性对的合取，才能检验已学规则向新属性对的迁移；只有孤立颜色标签时，无提示执行任意逻辑属于另一项未建立的问题。完整模式记忆、单属性、路由和联合模块方案应在同一留出集比较，再加入其他属性关系或指令顺序，检查一次成功是否依赖固定词序。模块连接与参数相加说明实现可行，留出组合上的正确结果才说明所测组合能力。

\subsection{全过程资源消耗的度量}
\label{sec:ket-evaluation-resources}

共享或迁移节省了适应费用以后，还需问取得和维持这些成绩共花了多少资源。来源训练与准备可能只发生一次，选择和适应按任务发生，服务、更新与恢复则随时间反复发生。沿用式~\eqref{eq:ket-reuse-amortization}，$C_{\mathrm{prep}}$计核算范围内的来源训练及准备，$C_{\mathrm{sel},k}$计所有候选评分和被淘汰试训，$C_{\mathrm{adapt},k}$只计此前未入账的最终适应。

设实际服务$m$个任务，观察区间为$[0,H]$。共享骨干一次前向同时给出三项输出，只发生一个推理事件；因此用实际事件集合$\mathcal E_{\mathrm{infer}}(H)$及事件增量成本$c_e$计推理，不能逐任务重复累加。任务$k$在区间内的更新费用为$C_{\mathrm{update},k}(H)$，再纳入共享维护与恢复，得到计算账目：
\begin{equation}
  \begin{aligned}
    C_{\mathrm{total}}(m,H)
    ={}&C_{\mathrm{prep}}+C_{\mathrm{shared}}(H)\\
       &+\sum_{k=1}^{m}\bigl(
        C_{\mathrm{sel},k}+C_{\mathrm{adapt},k}
        +C_{\mathrm{update},k}(H)\\
       &\hspace{24mm}+C_{\mathrm{recover},k}(H)\bigr)
        +\sum_{e\in\mathcal E_{\mathrm{infer}}(H)}c_e.
  \end{aligned}
  \label{eq:ket-evaluation-ledger}
\end{equation}
$C_{\mathrm{shared}}$仅包含未归入任务项的一次或周期性维护，$C_{\mathrm{recover},k}$包含失败重试与恢复。同一更新费用不能再进入共享项。若需要把推理事件分摊给任务，可另定非负份额$\rho_{e,k}$，对每个事件要求$\sum_k\rho_{e,k}=1$，于是各任务的$\sum_e\rho_{e,k}c_e$之和等于真实推理总费。分摊改变归因，不增加事件。

每项费用附事件或版本标识、时间、设备、用途、服务任务和复用条件；未获益的任务仍保留已发生费用。若选择阶段训练的检查点继续适应，适应项只加后续计算；若丢弃它后重新训练，则确实要登记第二次费用。这些记录让总和能逐项复查。

已有来源模型可以按本项目新增投入核算，也可以按含来源训练的生命周期核算，但比较表需使用一致口径，不能只免除某一方案的历史费。准备成本按实际任务数$m$与期限$H$摊销，设想的未来任务只能作情景分析。教师、回放记忆、生成器、特征缓存、索引和版本副本同样登记；教师调用已计入训练项时，不再另加同一次计算。

人工、设备计算与存储未必有合理的统一换算，应分别报告。标注与审核记人时，计算记设备型号、时长及必要运算量，内存记峰值，持久存储记权重、数据、缓存字节数与保留期；随时间累计的字节时也不同于峰值字节。恢复等待和服务停顿另记墙钟时间。参数少只能提示权重占用可能减少，不能代替延迟、吞吐和总计算的实测。

训练依据的制作与校验属于第~\ref{part:efficient-knowledge-use}部分，服务和维护的实现属于卷十二的系统主题。本章只用已经定义的账目、单位与去重规则比较投入，不因跨任务汇总而强行合并不同资源。

\section{评价数据与基准的选择}
\label{sec:ket-evaluation-data}

假如主张是能够进入新地区，随机拆分同一地区的照片即使得到高准确率，也没有检验这项变化。按类别分阶段可以控制学习次序，却不等于真实季节流。指标规定测什么，数据还必须让对应条件出现。评价基准提供可重复的任务与划分，其分数支持这些条件下的判断，不支持任意环境中的有效性。

\subsection{评价主张与数据覆盖的匹配}
\label{sec:ket-evaluation-coverage}

要检验花卉识别跨地区的变化，可先保持类别目标不变，再留出新的采集领域。DomainBed将多领域分类数据、算法与模型选择放入统一框架，领域由照片风格、来源或采集环境等因素区分。在一套数据内，变化通常落在领域而非类别目标。它区分训练领域验证、留一领域交叉验证及测试领域验证；最后一种用目标领域标签选配置，具有更强的先知权限，不能与无目标数据的领域泛化并列解释。这个框架检验的不仅是算法，也包括选择超参数的规则，但未由此检验新类学习、事实修正或长期维护\scite{gulrajani2021}

WILDS把医院、相机位置、地理位置及时间等自然变化接入应用任务的训练与分布外评价划分。有些设定分离训练和测试领域，有些领域重叠而群体比例改变；后一类需要关注弱势群体，不能都解释为整领域留出。与人为像素扰动相比，变化来自实际采集条件或群体组成，应用目标仍由子任务确定。报告须写明子数据集、领域或群体和官方指标，不能合成一个未定义的“WILDS分数”。这些自然场景也不覆盖全部未来偏移，跨医院分类成功并未提供跨任务少样本证据\scite{koh2021wilds}

若主张改为适应未见分类任务，就还需隔离任务或类别，单独留出照片不够。Meta-Dataset从多套图像数据抽取少样本任务，变化包括类别数、支持集规模及不平衡程度，并利用部分数据的类别层级。每个任务来自单一数据集；原协议将Traffic Signs与MSCOCO全部留作评价，其余来源另有类别划分。这使未见类别与整体未见数据来源可以分别检查。任务仍主要是图像分类，所得成绩不直接回答病斑分割、语言生成或任意新损失下的适应问题\scite{ket-triantafillou2019metadataset}

未见来源有时还改变成像方式。BSCD-FSL从miniImageNet基础类别准备模型，再在CropDiseases、EuroSAT、ISIC2018与ChestX上评价植物病害、卫星、皮肤病变和胸部X射线分类。研究控制基础模型，在相同抽样任务上比较元学习与普通微调，并采用5分类、每类5、20或50个支持样本。这些设定保持少样本分类形式，同时改变类别和成像领域；不同数据集的全部差异不能缩成单一、已控制的“领域距离”，更不能据此判断长期遗忘\scite{ket-guo2020-cross-domain-few-shot}

长期使用主张需要保留时间戳及反馈次序。CLEAR利用YFCC100M的时间信息构造2004至2014年的视觉变化序列，主要在固定类别空间中检验领域增量。除同一时间桶内的训练测试划分外，它还采用当前模型预测下一时间桶的流式协议：先封存预测，再获得标签，随后该桶可进入训练。这检验近未来预测，但这些数据进入训练后就不能继续充当独立历史保持测试。若同时研究遗忘，应另留历史审计集并登记额外资源。时间戳本身并未增加无限新类、无任务边界或事实撤回等条件。

连续预训练进一步涉及语料到达阶段与下游任务两种对象，应分别记录。每个语料阶段结束后，可以在固定下游任务上复测保持，再用隔离新任务测学习成本，同时计新增词元、旧语料回放及训练计算。新语料损失下降并不能说明下游能力保持；一次早期自监督预训练也不是长期持续预训练。CLEAR含有丰富无标签图像，但早期预训练的收益不能替代逐时间段更新后的这些检查。

对所在地这样的内容修改，数据应围绕事实依赖组织。MQuAKE-cf将事实链中部分内容改成反事实，MQuAKE-t考察真实时间变化；关系链其他部分保持不变，用多跳答案检验修改后果是否传播。它还包含按特定模型可回忆事实筛选的条件，需报告实体范围、链长、编辑数和问题生成方式。这种数据检验关联推理，而非仅检验编辑提示改写，也不是隐私删除或全部长期冲突的基准。

机器遗忘的数据则必须对应真实训练记录、撤回粒度和时机，并有匹配的重训练参照，例如SISA在Purchase、SVHN等任务中的删除实验。随机记录撤回改为整类或整用户撤回，可能同时改变数据规模与任务要求，应另列这些条件，不沿用原来的删除结论。

共享数据需让独立学习对照可信。花卉任务集应记录每张照片具备哪些标签，以及缺失是否与地区或病害严重度相关。Taskonomy在共同室内场景图像上提供多种标注，以固定输入像素研究来源任务向目标的迁移关系；这个关系分数并不直接等于联合训练收益。借其数据评价共享时，还需固定每任务标签、共同对象、划分和独立预算。若共享方案额外看到带物种标签的病害照片，就应分开增加监督与共享机制的作用\scite{ket-zamir2018-taskonomy}

组合数据需要真正留出完整关系，而不只是新的样本编号。SCAN把有限语法生成的导航指令转成动作序列，原语是基本动作或指令成分，其含义和组合规则固定。随机指令划分、长度外推以及某个原语只在单独指令中出现的划分，分别提出不同的组合要求。随机划分成绩不能覆盖后两类要求；受控语法有助于隔离新组合，却不包含真实视觉、开放语境和无限语法\scite{ket-lake2018-scan}

COGS以受控英语句子为输入，输出表达句子关系的逻辑形式，检验熟悉词语的新语法角色、结构重组及更深嵌套。它比SCAN涉及更多语言结构和语义角色，但仍由规则生成英语片段，未因此覆盖跨语言或视觉组合。应分列同分布测试与组合泛化集，并逐类型报告；只给均值可能隐藏词汇替换成功、结构重组失败的差异\scite{ket-kim2020-cogs}

\begin{samepage}
融合评价最后还要依据来源如何形成来选数据。Model soups比较同一预训练起点经不同微调配置得到的最佳单模型、权重平均与预测集成，其同目标任务收益不直接说明跨任务知识整合。多任务融合需另设逐来源原任务测试、联合请求及组合留出，并核对来源筛选和融合系数是否用过这些测试。平均后部署一个模型只说明推理时无需保存多个副本，候选训练与配方选择费用仍存在\scite{ket-wortsman2022-soups}

\end{samepage}

表~\ref{tab:ket-evaluation-claims}~用八项可多选标签整理上述评价对象，不将它们视为互斥类别。上块列出变化和必要留出，下块按相同行序列出基准、信息权限与不能推出的结论。同一实验若检验多项主张，就要逐行满足相应条件；“须声明的信息权限”列提示应报告的权限，不表示同名基准默认允许全部信息。

\begin{table*}[tbp]
  \centering
  \footnotesize
  \renewcommand{\arraystretch}{1.2}
  \begin{tabularx}{\linewidth}{@{}>{\raggedright\arraybackslash}p{24mm}
    *{2}{>{\raggedright\arraybackslash}X}@{}}
    \toprule
    \multicolumn{3}{@{}l@{}}{\textbf{(a) 评价主张、变化与留出}} \\
    \midrule
    对象标签（可多选） & 实际改变的条件 & 需要隔离的对象 \\
    \midrule
    \mbox{自然偏移} & 保持子任务目标，改变领域、采集条件或时期 &
      对应领域与独立对象组；时间主张另留整段时期 \\
    \addlinespace[2pt]
    \mbox{未见任务} & 变换类别、支持集规模或任务来源 &
      最终任务与类别；跨领域主张另隔离数据来源 \\
    \addlinespace[2pt]
    \mbox{真实时间流} & 按时间到达样本，标签空间可保持不变 &
      预测时尚未到达的时期；独立历史审计集 \\
    \addlinespace[2pt]
    \mbox{知识编辑} & 修改指定事实，并检验依赖后果 &
      未用于编辑的改写、关联题、无关题及历史题 \\
    \addlinespace[2pt]
    \mbox{机器遗忘} & 按既定粒度与时机撤回训练记录 &
      保留任务测试、撤回行为审计和攻击测试 \\
    \addlinespace[2pt]
    \mbox{多任务共享} & 改变任务搭配、标签覆盖或共享结构 &
      独立对象或场景组，并为各任务保留测试 \\
    \addlinespace[2pt]
    \mbox{组合泛化} & 已知成分进入新的关系、深度或次序 &
      完整组合及等价表达，同时保留基本成分覆盖 \\
    \addlinespace[2pt]
    \mbox{模型融合} & 选择来源参数、融合配方和服务请求分布 &
      各来源原任务测试、联合请求及组合留出 \\
    \bottomrule
  \end{tabularx}
  \par\medskip
  \begin{tabularx}{\linewidth}{@{}>{\raggedright\arraybackslash}p{24mm}
    >{\raggedright\arraybackslash}p{36mm}
    *{2}{>{\raggedright\arraybackslash}X}@{}}
    \toprule
    \multicolumn{4}{@{}l@{}}{\textbf{(b) 代表基准、信息与证据边界}} \\
    \midrule
    对象标签（可多选） & 数据或研究参照 & 须声明的信息权限 & 不能据此推出 \\
    \midrule
    \mbox{自然偏移} & \mbox{DomainBed}；\mbox{WILDS} &
      来源标签；是否使用目标选择数据 & 对任意未来偏移或新任务均可适应 \\
    \addlinespace[2pt]
    \mbox{未见任务} & \mbox{Meta-Dataset}；\mbox{BSCD-FSL} &
      元训练、元验证及目标支持证据 & 有限分类任务覆盖任意任务分布 \\
    \addlinespace[2pt]
    \mbox{真实时间流} & \mbox{CLEAR} &
      只读已到达信息，获得反馈后训练 & 近未来预测同时证明独立历史保持 \\
    \addlinespace[2pt]
    \mbox{知识编辑} & \mbox{MQuAKE} &
      编辑请求与获准使用的辅助事实 & 单跳正确说明全局一致或影响已删除 \\
    \addlinespace[2pt]
    \mbox{机器遗忘} & \mbox{SISA}中的\mbox{Purchase}、\mbox{SVHN}删除实验 &
      剩余数据、保存状态和重训练规则 & 拒答或攻击失败即验证严格删除 \\
    \addlinespace[2pt]
    \mbox{多任务共享} & 花卉任务集；\mbox{Taskonomy}数据 &
      每任务样本、标签及缺失情况 & 额外数据的收益全部来自共享机制 \\
    \addlinespace[2pt]
    \mbox{组合泛化} & \mbox{SCAN}；\mbox{COGS} &
      成分与规则可见，目标组合留出 & 符号组合成功涵盖任意多模态组合 \\
    \addlinespace[2pt]
    \mbox{模型融合} & \mbox{Model soups}；多来源任务测试 &
      来源模型、选择数据及任务头权限 & 参数可相加说明来源无损或新能力形成 \\
    \bottomrule
  \end{tabularx}
  \caption{主张与评价条件的对应表。两块按相同对象逐行连接：上块说明应改变和隔离什么，下块说明可参考的数据、须报告的权限及外推限制。标签可以同时选择，但各项条件不能互相代替。花卉任务集为教学设计，其余参照见正文引文；表中没有跨基准排名，删除实验也不等于形式化删除保证。}
  \label{tab:ket-evaluation-claims}
\end{table*}

\subsection{基准条件与信息来源的核查}
\label{sec:ket-evaluation-provenance}

“本次训练未用测试文件”仍不能回答预训练是否接触过测试信息。核查应区分相同输入、输入及标签、完整任务规则、完整组合和近重复记录。某物种的无标签照片与测试照片的物种答案提供不同信息；见过颜色、纹理词语也不等于见过“红色且有斑点”的完整指令。保留这些重叠类型，才知道哪项留出主张受到影响，不能仅勾选未经说明的“无污染”。

核查记录应包含来源语料或模型说明、时间范围、版本、可访问清单，以及去重对象和方法。哈希相等能发现完全相同文件，却不能排除图像近重复或文本改写，后二者需要相应检索与人工检查。报告阈值、重叠数、排除规则和漏检风险；只检查公开子集时，结论也只能是在该范围内未发现所查重叠，不能推广至未知语料。

预训练覆盖之外，还需登记本次运行的权限：旧原始数据、记忆与检查点能否访问，任务身份是否已知，标签何时到达，以及能否共同处理一批目标无标签输入。任务身份已知的分类与未知身份分类不具有相同条件。来源筛选、提示示例、阈值校准和人工失败分析读取的目标数据也属于开发访问，即使没有参与梯度计算。

一条结果应能追溯到具体访问条件，例如新地区使用了30张带标签选择照片、80张适应照片及未标注校准批次，而未用最终测试标签选择来源。“标准协议”这类称谓无法替代记录。第~\ref{sec:ket-evaluation-splits}~节将这些权限落实为整体留出与共享规则；查到影响主张的重叠后，应重做独立评价或收窄结论，不能只附脚注而保留原强度。

\subsection{补充场景与失败样本的构造}
\label{sec:ket-evaluation-stress}

若花卉系统要跨季节长期服务，而现有基准只测一次跨地区迁移，就还需补充稀有病害、遮挡、旧地区回归及更长序列。持续编辑则需考察同一事实再次更新、冲突请求与依赖前提缺失。补测条件由目标中尚未覆盖的要求决定，纳入规则在看结果前固定。事后发现的失败可以整理为诊断集，但不再代表预先未定义总体的失败率。

为了区分失败原因，可以构造成对条件。组合比较保持颜色、纹理、对象数量和背景分布，改变属性关系是否在训练中共同出现；分布适应固定标签与样本数，改变光照；关联编辑固定事实、模板和编辑数，改变依赖链长。若更长链也含更多冷门实体，就同时改变了熟悉度，不能把全部差异归给推理深度，应再匹配实体条件。

成对条件也必须具有有效语义。迁址请求缺少生效日期时，历史和当前问题可能没有唯一答案，应补日期或把明确报告信息不足列为目标。冲突请求须给出版本优先级或待核验状态，不能让模型猜哪条可信。合成病害若不符合成像及病理条件，就只能评价所定义图案的响应，而非临床或田间病害识别。

本章矩阵用于复算定义，添加误差棒也不会使人工成绩成为实验。受控实验用于区分候选解释，真实数据用于检验部署条件，三类材料应各注明用途。正式补充数据还需保存生成规则、种子、采样范围、质量检查及完整纳入清单，让所有方法接受同样条件，并保留成功、失败与无法评价的记录。

\section{对照实验与评价协议的设计}
\label{sec:ket-evaluation-protocol}

即使使用同一套照片，预测前读到标签与预测后才读到标签仍会产生不同问题。因此数据确定后，还要固定访问时机、比较方案、资源条件与重复所覆盖的变化。若中途修改选择规则或成功阈值，新的分析应标为探索性；用来确认新主张的数据需要重新独立安排。

\subsection{数据划分与访问协议的设计}
\label{sec:ket-evaluation-splits}

要让花卉数据检验预定变化，应按相应对象整体留出。样本层隔离照片，并将同植株连拍、裁剪与增强归为同一对象组；类别层隔离最终物种；任务层隔离未参与元训练及选择的预测要求；领域层隔离地区、设备或来源；时间层只允许过去信息预测未来。照片不同仍可能是同一植株，物种不同仍可能来自同一区域，支持与查询照片分开也可能仍属于已见任务族。这些留出条件分别成立，不能互相代替。

例如在各地区、各物种内部随机拆照片，评价范围仍是已见类别与领域中的新样本。将丙、丁物种整体留下才增加未见类别要求，但若照片来自已见地区，仍未检验新领域。再留出地区乙、以2026年7月分隔过去与未来，分别增加领域和时间条件。多个层级联合隔离可能减少可用数据，报告应直接说明控制了哪些层级，不用“严格划分”代替它们。

图~\ref{fig:ket-evaluation-boundaries}~将选择、适应与最终计分分开。来源数据形成候选，目标选择集决定来源、超参数及停止规则，适应集用于拟合，最终评测集只在方案锁定后计分。少样本规则可在开发任务上选择，最终任务仅使用获准支持证据。若预定协议允许选择集并入适应集，应标注两用途共享、按唯一样本核算标签，不能再把该集视为独立证据；最终评测仍在外部。

底部三个网格复用同一组九张照片，只改变隔离方式。样本留出后，训练仍含A、B、C物种及两个地区；留下C整行才要求未见物种识别，留下第3列才要求在未见地区乙识别已见物种。图中变化的不是样本池，而是测试必须跨越的边界，这解释了为什么无重复照片仍不足以支持类别或领域外推。

元学习中还要区分任务内与任务间的反馈。沿用第~\ref{sec:ket-transfer-task-splits}~节，$\mathcal T_k\sim\Pi$表示从任务分布抽取学习问题，$S_k$用于支持适应，$Q_k$用于查询。元训练的查询标签更新学习规则，所以不是独立最终证据；元验证比较候选规则，选定并锁定后，交给最终任务的是规则，而非验证任务及其标签。最终$S_k$提供获准适应证据，$Q_k$标签只计分。

新抽一组支持照片不等于留下新类别或领域，任务抽样仍须满足对应层级隔离。查询输入权限也应声明：直推式协议（transductive protocol）允许共同利用同一任务的全部查询输入，归纳式协议（inductive protocol）逐个接收查询。两者没有相同信息条件，不能只因支持标签数一致就混合比较。

\begin{figure*}[htbp]
  \centering
  % Solid: fitting/update or labeled rule handoff; dashed: selection evidence.
% Dotted: scoring. Task collections never flow into one another.
\begin{tikzpicture}[
  x=1mm,y=1mm,font=\figurefont\footnotesize,text=BookInk,
  box/.style={draw=BookMuted,line width=0.7pt,minimum height=13mm,
    text width=27mm,align=center,inner sep=1.5mm},
  train/.style={-{Latex[length=2mm]},draw=BookTeal,line width=0.9pt},
  select/.style={-{Latex[length=2mm]},draw=BookAmber,line width=0.8pt,dashed},
  score/.style={-{Latex[length=2mm]},draw=BookBlue,line width=0.8pt,dotted},
  title/.style={anchor=west,inner sep=0pt,font=\figurefont\small\bfseries}
]
  \path[use as bounding box] (0,-178) rectangle (169,17);
  \node[title] at (0,13) {(a) 开发与独立最终评价};
  \node[box,fill=FigureInput!10] (source) at (16,0) {来源训练样本\\形成候选};
  \node[box,fill=FigureModel!10] (knowledge) at (60,0) {候选模型\\确定复用配置};
  \node[box,fill=FigureModel!10] (adapt) at (105,0) {目标适应\\执行已选规则};
  \node[box,fill=FigureOutput!10] (frozen) at (149,0) {锁定模型\\封存预测};
  \draw[train] (source.east) -- (knowledge.west);
  \draw[train] (knowledge.east) -- (adapt.west);
  \draw[train] (adapt.east) -- (frozen.west);

  \node[box,fill=FigureInput!10] (selection) at (60,-26)
    {目标验证样本\\来源、训练配置、停止};
  \node[box,fill=FigureInput!10] (support) at (105,-26)
    {目标训练样本\\拟合任务模型};
  \node[box,fill=FigureInput!10] (test) at (149,-26)
    {最终测试样本\\只读计分};
  \draw[select] (selection.north) -- (knowledge.south);
  \draw[train] (support.north) -- (adapt.south);
  \draw[score] (frozen.south) -- (test.north);
  \draw[draw=BookMuted,line width=0.6pt]
    (selection.south) -- (60,-38) -- (105,-38) -- (support.south);
  \node[anchor=north,inner sep=1mm] at (82.5,-38)
    {仅按预定协议共享};
  \node[align=left,anchor=west,inner sep=0pt,text width=30mm] at (0,-27)
    {最终结果不回流\\选择与适应};

  \node[title] at (0,-54) {(b) 元学习：任务层留出};
  \node[box,text width=31mm,fill=FigureInput!10] (metatrain) at (18,-68)
    {元训练任务\\$S_k$拟合，$Q_k$更新};
  \node[box,text width=31mm,fill=FigureModel!10] (candidates) at (62,-68)
    {候选规则};
  \node[box,text width=31mm,fill=FigureInput!10] (metaval) at (18,-88)
    {元验证任务\\试用候选规则};
  \node[box,text width=31mm,fill=FigureModel!10] (lockedrule) at (62,-88)
    {锁定规则};
  \node[box,text width=31mm,fill=FigureOutput!10] (metatest) at (62,-111)
    {最终测试任务\\$S_k$适应，$Q_k$计分};
  \draw[train] (metatrain.east) --
    node[midway,above,inner sep=1mm] {学习} (candidates.west);
  \draw[select] (candidates.south) --
    node[midway,right,inner sep=1mm] {选择} (lockedrule.north);
  \draw[select] (metaval.east) --
    node[midway,above,inner sep=1mm] {验证} (lockedrule.west);
  \draw[train] (lockedrule.south) --
    node[midway,right,inner sep=1mm] {交接规则} (metatest.north);
  \node[anchor=west,align=left,text width=33mm,inner sep=0pt] at (0,-111)
    {任务集合分开\\未见类别另做\\类别层留出};

  \node[title] at (91,-54) {(c) 连续学习：时间顺序};
  \node[box,text width=56mm,minimum height=9mm,fill=FigureOutput!10]
    (predict) at (128,-68) {状态$\sigma^{(t)}$：先预测并封存};
  \node[box,text width=56mm,minimum height=9mm,fill=FigureInput!10]
    (feedback) at (128,-85) {反馈到达：计分并按协议更新};
  \node[box,text width=56mm,minimum height=9mm,fill=FigureModel!10]
    (next) at (128,-102) {状态$\sigma^{(t+1)}$：仅服务未来};
  \draw[score] (predict.south) -- (feedback.north);
  \draw[train] (feedback.south) -- (next.north);
  \node[anchor=west,inner sep=0pt] at (91,-115)
    {独立运行重置；流内状态按协议保留};

  \node[title] at (0,-132) {(d) 同一样本池，不同留出层级};
  \foreach \shift/\kind/\heading in {0/1/样本留出,55/2/类别留出,110/3/领域留出} {
    \node[anchor=west,inner sep=0pt] at ({\shift+10},-140) {\heading};
    \foreach \r/\letter in {1/A,2/B,3/C} {
      \foreach \c in {1,2,3} {
        \ifnum\kind=1
          \ifnum\r=\c
            \fill[BookAmber!18]
              ({\shift+10+7*(\c-1)},{-146-7*(\r-1)})
              rectangle ({\shift+17+7*(\c-1)},{-153-7*(\r-1)});
          \fi
        \fi
        \ifnum\kind=2
          \ifnum\r=3
            \fill[BookAmber!18]
              ({\shift+10+7*(\c-1)},{-146-7*(\r-1)})
              rectangle ({\shift+17+7*(\c-1)},{-153-7*(\r-1)});
          \fi
        \fi
        \ifnum\kind=3
          \ifnum\c=3
            \fill[BookAmber!18]
              ({\shift+10+7*(\c-1)},{-146-7*(\r-1)})
              rectangle ({\shift+17+7*(\c-1)},{-153-7*(\r-1)});
          \fi
        \fi
        \draw[BookMuted,line width=0.5pt]
          ({\shift+10+7*(\c-1)},{-146-7*(\r-1)})
          rectangle ({\shift+17+7*(\c-1)},{-153-7*(\r-1)});
        \node[inner sep=0pt] at ({\shift+13.5+7*(\c-1)},{-149.5-7*(\r-1)})
          {\letter\c};
      }
    }
  }
  \node[anchor=west,inner sep=0pt] at (0,-174)
    {A、B、C为物种；第3列来自地区乙，其余来自地区甲；阴影为留出。};
\end{tikzpicture}

  \caption{数据用途、任务留出与时间访问边界。实线为拟合、协议更新或规则交接，虚线为选择依据，点线为只读计分；箭头不表示任务集相互并入。最终任务只接收锁定规则，查询标签不返回本轮开发；选择与适应共享须预先声明。下方九张教学照片保持相同，阴影分别标出样本、类别及领域留出。在线预测先封存，反馈只改变未来状态，不重算过去成绩。}
  \label{fig:ket-evaluation-boundaries}
\end{figure*}
\FloatBarrier

如果测试输入还会改变模型，隔离的不仅是数据，也包括状态。以$\sigma^{(t)}$表示参数、归一化统计、优化器、记忆、缓存及随机数等可变状态；独立任务或运行开始时恢复约定起点。同任务内跨样本是否保留、结束后是否丢弃，都需固定。只恢复权重却沿用上一任务统计或缓存，并未完成要求的重置。

在线评价的\term{先预测后更新}（predict-then-update）规则用当前状态生成并封存$\hat y_i$，待标签到达才计分与更新，新的状态服务以后样本。若第$i$个标签延迟$d_i$步，预测时只能读到此前已到达反馈，不能提前使用$y_i$。测试时适应也可另定先用当前无标签输入更新再预测的协议，但应列明无标签样本范围及是否整批访问。

读到真实标签、更新后再对同一对象计分，得到的是训练拟合而非原来的预测成绩。预定在线规则允许流式标签训练未来状态，不允许事后替换已经封存的预测。

持续序列可以来自真实时间、受控排列或旧环境回归，三者回答不同问题。真实时间必须保持因果访问，不能为改善均值任意打乱；受控排列可重复抽取顺序来检验顺序效应；环境回归问再次遇到旧条件时能否恢复与继续学习。应分列结果；若另设独立最终审计，审计记录不进入在线反馈流。

编辑划分则按请求及关系族隔离：原请求用于修改，改写、跨语言、关联题和无关题分别留出。重复直接答案的模板不构成多个独立事实成功；多跳测试可以提供被修改事实，不能同时把最终多跳答案作为示例。若要研究权限本身，可固定模型和数据，分别设置有无源数据、任务身份或测试状态的组。它们测的是权限变化，强权限成绩不能称为同权限下的算法收益。

\subsection{比较基线与消融实验的设计}
\label{sec:ket-evaluation-baselines}

来源初始化取得更高成绩，可能因为来源有帮助，也可能因为从头训练没有得到合理优化。基线应针对这种直接替代解释设置：保持结构与目标数据，比较从头训练、普通微调及冻结来源或随机表示；持续保持与朴素顺序训练、同记忆预算简单回放比较；共享与各任务独立学习及同总预算独立方案比较。各基线都需合理调参，不能仅用未收敛的随机模型制造优势。

检验专用梯度协调器时，单任务对照还取消了共享结构，无法单独隔离协调器。应固定骨干、任务头、数据、损失单位、起点、基础优化器和停止规则，设置等权联合训练，并在预定预算内搜索固定损失权重或采样概率。候选范围、验证选择与搜索预算预先列出，测试只评价锁定方案；逐任务梯度、协调求解和淘汰配置均记账，结果保留逐任务取舍。Xin等在所测语言任务与方法中发现，专用方法未超出简单采样搜索的性能取舍，因此充分调参的简单联合训练是必要的直接对照。这一观察不覆盖全部协调器，也不覆盖未纳入该研究的Nash-MTL\scite{ket-xin2022-mto}

融合至少与选择集确定的最佳单来源、直接预测集成比较。按最终测试逐任务挑最佳来源属于先知权限，应明确标注；多头集成的任务身份优势也不能隐去。融合若效果低于集成而推理成本更低，可以是一种取舍，但需并列报告效果与资源，不能只用较弱单模型声称全面占优。模块组合还应与单来源路由、移除模块及无信息输入替换比较，检查各成分是否真正影响所需结果。

元学习比较需保持预训练来源、骨干容量和目标支持集，再设置普通适应、仅用初始化不适应以及完整元学习适应。普通适应缺少相同预训练时，学习规则就不是唯一差异。来源选择则可比较默认、随机、代理评分与等预算目标试训，连同搜索开销报告。最终测试上挑出的最佳来源只能作先知诊断，不能作为可部署选择器的成绩。

直接基线之后，还可通过部件干预检查机制解释。\term{消融实验}（ablation study）移除或替换某个部件，再观察完整方案如何改变。移除回放可能同时减少旧数据权限，移除共享可能改变参数与计算，移除适应可能减少更新次数；这些伴随变化都需登记。还可将节省的适应预算用于普通更新，区分专门机制与更多更新。消融并不天然只改变一个因素。

允许全部历史数据的联合训练可提示受限记忆的代价，但其权限强于许多持续协议，应另列为强参照而非同条件排名。有限网络、计算与优化得到的分数也不是数学上界，其他优化过程仍可能超过它。因而报告应称“本配置下的全历史联合训练参考”，不称最优保证。

\subsection{资源预算与重复实验的设计}
\label{sec:ket-evaluation-repeats}

假如整项开发只有固定预算，来源选择、失败试训与最终适应都必须在其中支付；若目标是同一效果，则记录完整过线过程和未过线运行。来源准备范围、共享摊销任务数、服务请求量与截止时刻应预先固定。方案在预算内无法完成时，记录不可执行或未完成，不仅筛出便宜的成功运行。

“资源相同”要落实到具体量。相同标签数可能涉及不同标注工时，相同步数可能使用不同批量与网络，相同GPU小时可能来自不同设备和利用率，相同权重大小也不约束内存峰值。可以分别设置标签、实际计算或内存约束，不强求所有量同时相等。未匹配的差异仍在结果表中保留，并限制相应归因。

重复之后要计算区间，还需先说明平均收益针对哪种随机变化。借用第~\ref{sec:ket-evaluation-target}~节的$(6,6,4)$个百分点数值，另设仅演示计算的教学协议：固定任务、训练与支持数据、预先封存测试照片、配置和超参数，只重新抽取运行随机性。每对方案共用可匹配条件，不同对独立抽种子，故独立单位是完整的一对运行，而不是其中的照片。这是使用相同数值说明新协议，并未把原验证记录改称独立测试。

令$d_r$为第$r$对的准确率差，$\mu_d=\mathbb E_{\mathrm{run}}[d_r]$是在上述固定条件下对完整运行随机性取平均的收益；$\mathbb E_{\mathrm{run}}$不对新任务或新照片取期望。为说明小样本$t$区间的精确推导，先引入一个连续运行差的理想模型，假设其中的差值独立同分布且为正态。固定100张照片时，准确率差只能取有限格点值，并不精确服从连续正态分布；将本例的三个人工差值代入$t$公式只能演示近似计算，不能由原计分设定获得有限样本精确覆盖。式~\eqref{eq:ket-evaluation-run-summary}~的$s_d$描述单次差值波动，而$\widehat{\operatorname{SE}}(\bar d)=s_d/\sqrt R$估计有限运行数下的均值标准误。

在这个连续正态理想模型下，均值误差经估计标准误标准化后服从自由度$R-1$的Student $t$分布。预定95\%覆盖率，将两侧各留2.5\%，以$t_{\nu,p}$表示自由度$\nu$的$p$分位数，得到：
\begin{equation}
  \bar d\ \pm\ t_{R-1,0.975}\,
    \frac{s_d}{\sqrt R}.
  \label{eq:ket-evaluation-paired-interval}
\end{equation}
将三次教学差值作为近似演算输入时，$R=3$、自由度为2，计算为：
\[
  \begin{aligned}
    \bar d&=\frac{16}{3},\qquad s_d=\sqrt{\frac43},\\
    \widehat{\operatorname{SE}}(\bar d)&=\frac{\sqrt{4/3}}{\sqrt3}
      =\frac23,\\
    \bar d\pm t_{2,0.975}\frac23
      &\approx5.333\pm4.303\times0.667\\
      &\approx[2.46,\,8.20]\text{个百分点}.
  \end{aligned}
\]
95\%指构造规则的长期覆盖率：在连续正态理想模型成立时，反复抽取成组运行，所得区间约有95\%覆盖固定的$\mu_d$，不表示某次运行有95\%概率获益。对固定照片上的离散准确率差，上述数值区间只是近似示范，不继承这项精确覆盖结论；三个人工数值也不能检验近似是否合适，更不是实际重复次数建议。区间只覆盖运行随机性，不含新照片、新任务、支持集重抽或重新调参的不确定性。基准研究应分别登记这些随机来源；均值差与胜出概率是不同统计量，不能借用后者公式来计算前者区间\scite{ket-bouthillier2021-variance}

若评价多个任务，就还需区分同任务重跑与更换任务。以$g_{k,r}$表示任务$k$的第$r$次配对收益，设各任务重复$R$次，$K$个任务等权。先在任务内求$\bar g_k$，再求任务等权均值$\bar g$；分别用任务内离差和任务均值离差描述两种变化：
\begin{equation}
  \begin{aligned}
    \bar g_k&=\frac1R\sum_{r=1}^{R}g_{k,r},&
    \bar g&=\frac1K\sum_{k=1}^{K}\bar g_k,\\
    s_{\mathrm{run}}^2
      &=\frac1K\sum_{k=1}^{K}
        \frac{\sum_{r=1}^{R}(g_{k,r}-\bar g_k)^2}{R-1},\\
    s_{\mathrm{task}}^2
      &=\frac{\sum_{k=1}^{K}(\bar g_k-\bar g)^2}{K-1}.
  \end{aligned}
  \label{eq:ket-evaluation-two-variations}
\end{equation}
这里要求$R,K>1$，各分母对应其样本方差。$s_{\mathrm{run}}^2$平均同任务内运行波动，$s_{\mathrm{task}}^2$描述任务均值之间的差异，后者仍含有限重复带来的估计噪声，不是未经建模就分离出的任务总体方差。两者都是描述量，不是$\bar g$的标准误，不能直接作均值区间。

例如两个任务分别记录$(5,6,7)$与$(-3,-2,-1)$个百分点，任务均值为6与$-2$，总体等权均值为2。两组任务内方差都为1，故$s_{\mathrm{run}}=1$个百分点；均值的离差平方和为32，故$s_{\mathrm{task}}=\sqrt{32}\approx5.66$个百分点。更多重跑可减少任务均值的估计波动，却没有增加独立任务覆盖。正2的总均值也没有消除第二任务的负收益。

如果不采用正态运行差假设，可以依据已有独立单位近似抽样分布。\term{配对自助法}（paired bootstrap）有放回重抽这些单位，并让两方案始终使用相同索引，以保存对应关系，再观察统计量的波动。对固定任务和测试集的运行均值，可执行以下步骤。
\begin{enumerate}
  \item 先保存两方案每对分数与共同随机条件，固定统计量、任务权重、95\%置信水平、重抽次数$L$及种子。例如$L=10000$只表示重抽计算一万次，不是新增一万次训练实验。
  \item 每轮从$R$对运行中有放回抽$R$个索引，两方案使用同一串索引，逐对求差再平均。若汇总多个固定任务，沿用预定权重，不让重复更多的任务取得额外权重。
  \item 完成$L$轮后，取均值差的经验2.5\%和97.5\%分位数作为百分位自助区间，并报告原运行数、抽样单位、分位数规则与$L$。其覆盖率是近似的；三对记录无论重抽多少次，也不能补足原分布信息。
\end{enumerate}

多任务抽样要保留共同模型造成的依赖。一个元训练种子产生的模型若服务全部$K$个任务，抽中运行$r$时应整组带入$(g_{1,r},\ldots,g_{K,r})$，不能逐任务重抽不同种子。各任务确实独立训练时，才可分别重抽运行。若问题改为固定模型面对新测试对象的收益，应重抽独立对象组，两方案共用组索引并保留组内照片；这种区间覆盖对象抽样，不覆盖重训练。

从固定任务转向$\Pi$中的新任务，需要任务抽样本身有依据。若任务确由$\Pi$独立同分布抽取，且只评价一对固定元模型，可重抽$K$个任务索引，带入每项完整支持、查询和结果，重算任务均值。若还独立重复元训练，就有任务与种子两个交叉层级：分别抽任务和种子组，所有任务共用抽中的种子串，保留对应交叉记录；这是在任务抽样与种子抽样按设计独立时的一种分组近似。

支持集或查询对象也被重抽时，仍按实际设计保留其组内依赖，不能将任务、种子和照片摊平成独立样本。人工选出的两个任务不会因重抽而代表广泛总体。候选搜索随机性也只有重复整个选择过程才被覆盖，固定已经选出的配置不能提供这份变化。

逐项95\%覆盖并不是全部任务同时95\%覆盖。若事先决定共同判断$K$个任务，令总漏覆盖概率$\alpha=0.05$，Bonferroni校正（Bonferroni correction）为每项分配$\alpha/K$，即构造$1-\alpha/K$的双侧区间。每项满足上述连续正态理想模型时，式~\eqref{eq:ket-evaluation-paired-interval}~改用$t_{R-1,\,1-\alpha/(2K)}$。由各项漏覆盖概率之和至多$\alpha$，可得全部区间同时覆盖至少$1-\alpha$，不要求任务独立，代价是区间较宽。若把$t$公式近似用于固定照片的离散准确率差，校正后的边际与联合覆盖也仍是近似的。

自助法可相应取$\alpha/(2K)$及$1-\alpha/(2K)$分位数，但边际覆盖仍为近似，不能继承有限样本精确保证。这只涉及预定的$K$项；同时比较方法、预算或事后选任务时，应扩大比较族或另设确认数据。

顺序随机性也须单独检验。固定排列换种子只能描述该顺序下的波动；受控实验应报告抽取哪些排列及是否包含旧环境回归。真实时间不能倒放来增加重复，可以在因果边界内重复初始化，再用其他时间窗口或实体群检查外推。攻击与删除沿用可信性部分的验证口径，维护费用沿用第~\ref{sec:ket-evaluation-resources}~节和式~\eqref{eq:ket-evaluation-ledger}，不能在汇总时改成普通准确率或参数量。

\section{评价结果与证据边界的分析}
\label{sec:ket-evaluation-analysis}

得到完整分数后，还要回答开篇的三个疑问：谁因更新受益或受损，改善由哪些因素支持，已经测试的结果能推广到哪里。收益分布、机制比较与范围判断承担不同职责，都需回到原主张；记录齐全或平均较高本身并没有完成这些判断。

\subsection{收益分布与能力取舍的分析}
\label{sec:ket-evaluation-distribution}

固定预算决策先读整条曲线，再看预定终点。图~\ref{fig:ket-reuse-learning-curves}~的迁移优势从20标签时17个百分点缩至160标签时1个百分点，在已测点仍为正但幅度不同；第~\ref{sec:ket-reuse-prediction}~节的另一组构造则由早期领先转为后期落后。单个早期点因此不能代表预定后期检查点的结果。长序列还需并列旧任务成绩与后续学习成本，才能看到保持分数未变时的学习能力下降。

同一批结果换用平均单位，还可能改变改善方向。设地区甲900张照片的准确率由90\%升至92\%，地区乙100张由80\%降至70\%。教学构造的样本平均由$(810+80)/1000=89\%$升至$(828+70)/1000=89.8\%$，增加0.8个百分点；地区等权平均却由85\%降至81\%，下降4个百分点。前者按当前流量评价总正确数，后者让两地区同权；最差地区降至70\%又揭示弱项。稀有类、关键任务与预定权重结果应并列，不事后挑选有利均值。

同效果下的成本排序还依赖实际服务规模。图~\ref{fig:ket-evaluation-utility}~沿用第~\ref{chap:ket-reuse-selection}~章账目：复用先花准备费，逐任务费用却较低；从头训练没有这份准备费，随任务数增加更快。交点前后较便宜的方案反转，说明单次使用与多次使用不应混评。实际分析还需将式~\eqref{eq:ket-evaluation-ledger}~的维护、恢复和服务费用补回；若要合成效果与计算的效用，另说明权重及换算依据。

\begin{figure*}[htbp]
  \centering
  % Teaching costs from the reuse-selection chapter, equal target performance.
% C_reuse(m) = 120 + 10*m; C_scratch(m) = 30*m; all m = 1,...,10.
% Coordinate map: x = 24 + 12*m mm; y = 16 + 0.16*C mm.
\begin{tikzpicture}[
  x=1mm,y=1mm,font=\figurefont\footnotesize,text=BookInk,
  axis/.style={draw=BookInk,line width=0.7pt,-{Latex[length=1.8mm]}}
]
  \path[use as bounding box] (0,0) rectangle (169,82);
  \draw[BookTeal,line width=1pt] (29,76) -- (41,76);
  \fill[BookTeal] (35,76) circle (1mm);
  \node[anchor=west,inner sep=1mm] at (42,76) {复用};
  \draw[BookBlue,line width=1pt,dashed] (81,76) -- (93,76);
  \draw[BookBlue,fill=white,line width=0.8pt] (86,75) rectangle (88,77);
  \node[anchor=west,inner sep=1mm] at (94,76) {从头训练};

  \foreach \v in {0,100,200,300} {
    \draw[BookRule,line width=0.5pt] (24,{16+0.16*\v}) -- (144,{16+0.16*\v});
    \node[anchor=east,inner sep=1mm] at (23,{16+0.16*\v}) {$\v$};
  }
  \draw[axis] (24,16) -- (24,69);
  \draw[axis] (24,16) -- (151,16);
  \foreach \m in {0,1,2,4,6,8,10} {
    \draw[BookInk,line width=0.6pt] ({24+12*\m},16) -- ({24+12*\m},14.5);
    \node[anchor=north,inner sep=1mm] at ({24+12*\m},14.5) {$\m$};
  }
  \node[rotate=90] at (5,42) {总计算成本 $C$（GPU小时）};
  \node at (86,3) {实际服务任务数 $m$};

  \draw[BookTeal,line width=1pt]
    plot[domain=1:10,samples=10] ({24+12*\x},{16+0.16*(120+10*\x)});
  \draw[BookBlue,line width=1pt,dashed]
    plot[domain=1:10,samples=10] ({24+12*\x},{16+0.16*(30*\x)});
  \foreach \m in {1,...,10} {
    \fill[BookTeal] ({24+12*\m},{16+0.16*(120+10*\m)}) circle (1mm);
    \draw[BookBlue,fill=white,line width=0.8pt]
      ({23+12*\m},{15+0.16*(30*\m)})
      rectangle ({25+12*\m},{17+0.16*(30*\m)});
  }
  \draw[BookMuted,line width=0.7pt,dotted] (96,16) -- (96,44.8);
  \draw[BookMuted,line width=0.7pt] (96,46) -- (96,57);
  \node[anchor=south,inner sep=1mm] at (96,57) {$m=6$持平};
  \node[anchor=north west,inner sep=0pt,text=BookMuted] at (31,68)
    {成本越低越好};
\end{tikzpicture}

  \caption{同效果条件下，实际服务任务数$m$改变费用排序的教学构造。复用为$C_{\mathrm{reuse}}=120+10m$，从头训练为$C_0=30m$，单位GPU小时；相同推理费相消，维护等费用暂略。两者在$m=6$均为180，此后复用较低。点对应整数任务数，连线只引导观察，不表示额外测量。}
  \label{fig:ket-evaluation-utility}
\end{figure*}

连续编辑沿第~\ref{sec:ket-evaluation-editing}~节四组问题跟踪错误，整合则并列来源退化与组合留出失败。若目标间出现取舍，依据第~\ref{chap:ket-sharing}~章预定的主辅或共同交付要求判断，同时列出受益、受损及方向不确定的对象。总体收益不改变这些约定。

\subsection{机制贡献与替代解释的分析}
\label{sec:ket-evaluation-mechanisms}

假如共享方案比独立方案准确，但也多看了标注照片，当前比较还不能把改善全归给共享。更强预训练、额外监督、更多训练计算以及融合后增加的容量都可能构成替代解释。将完整方法、强基线和关键消融放入同一信息与预算表，检查哪些条件确已固定，哪些差异仍未分离，再限定机制结论。

来源代理评分还需与真实适应收益分别比较。固定适应预算后，高分来源是否领先、试训是否改变排名、被淘汰者是否仅在早期收敛较慢，都可检验。若代理排名与最终排名呈不完全的正秩相关，只能说明相对次序存在一定的一致倾向，仍可能出现逆序对或选错首位。即使两者次序完全一致，全部候选仍可能输给无迁移对照。试训得到更高终点成绩也可能花更多标签和计算，应按第~\ref{sec:ket-evaluation-target}~节分别评价固定预算收益与达标成本。

保持约束的收益也可能伴随后续适应变慢。固定记忆、新数据和骨干，只改变约束强度，将旧任务成绩、新任务曲线与总计算并列；若同时增加回放或更换骨干，就无法只归因于约束。消融直接支持的是所实施干预对系统的影响，只有伴随差异得到控制，才更接近目标因素的机制解释。

“梯度较一致的任务共享收益较高”若仅来自观察，仍是相关关系，不能说一致性导致收益。固定其他条件改变任务搭配，可以检验更具体的干预假说；在新地区实际工作负载上复现，才增加部署条件的证据。近期预印本同样需要公开配置、重复、消融和跨数据复核，发表时间及框架名称不替代检查。只在少数模型上验证，就保留这个范围。

\subsection{适用范围与证据缺口的分析}
\label{sec:ket-evaluation-limits}

一条少样本结论可以写为：在已测类别配置的图像分类任务中，提供类别描述，并允许每个任务使用五个带标签支持样本；在这一固定预算下，相对使用相同预训练的普通适应，目标查询集准确率更高。它不能缩写成无需监督的迁移，因为支持标签、任务范围、起点与预算都承担限定。若来源选择仅凭无标签代理分数，就还未直接确认目标收益；相同无标签输入仍可能对应不同标签规律，增加同类评分不能补出识别依据。相应信息边界可参考\BookChapterRef{chap:transfer-learning-theory}{迁移学习的可学习性}。

持续保持与后续可学性限于实际序列长度、到达规则和容量。三阶段教学矩阵用于解释指标，不能证明无限任务增长；固定类别的真实时间流也没有检验开放世界中新类不断加入。已测事实链一致不表示全部隐含依赖更新，符号指令组合成功也不覆盖未测跨模态或跨语言组合。第~\ref{chap:ket-overview}~章连接了这些研究问题，但概念联系不能代替新范围的数据与对照。

未检出失败、在指定测试上达标、在明确假设下获得保证，具有不同证据强度。有限攻击未发现泄漏可能受攻击能力限制；未检出收益既可能来自差异小，也可能来自任务或运行不足。统计可检测的极小差异也未必满足实际使用需要。报告需并列差异幅度、不确定性、成功要求与未测范围，不能只给显著或不显著。

证据不够时，应收窄到已支持条件，或按缺口安排新评价。领域不足需要新的领域留出，费用遗漏需要完整账目，关联失败需要成对关系链与明确前提；拼接不同权限、任务和分母的高分不能补足其中任一项。最终回到开篇，分别判断新地区整体改善是否伴随关键任务受损、旧能力相对什么时刻保持、组合是否超出来源选择，并给出取得和维持这些结果的总成本。

\section{本章小结}
\label{sec:ket-evaluation-summary}

判断一次演进或迁移是否值得采用，需要让主张、指标、数据和协议对应同一问题。指标给出对象、参照与分母，留出让所声称的变化实际接受测试，对照及权限记录检查条件差异，生命周期账目记录获取和维持收益的代价。固定预算更准确、达标少用标签、保持旧能力、纠正事实与处理新组合各有证据要求；它们可以同时成立，也可能发生取舍，一个综合分数不能替代其中任何一项。

三任务例子解释了同一成绩为何既低于历史峰值又高于刚学完：遗忘与BWT采用不同参照，学前收益与适应成本又描述不同阶段。合法事实更新按有效时间验证，不与有害遗忘混淆；撤回训练影响则另以未使用撤回记录的学习过程为参照。共享和融合仍需保留任务或来源退化，组合主张还需让基本成分可见、完整组合留出，不能由参数可连接推断。

可用结论应同时列出获益与受损对象、未过线运行、选择和适应所用信息，以及任务、时间和预算限制。平均改善只有放在这些条件中，才支持具体复用决策。若数据未覆盖关键要求、收益方向仍不确定或全程费用无法接受，评价同样提供暂不继续复用的依据，而不是要求每次实验都给出成功结论。
